Let $Q$ be a nontrivial unital additively idempotent semiring such that
$qq'=0$ implies $q=0$ or $q'=0$. Such a semiring is called \emph{entire}.
The map
\[
 \chi:Q\longrightarrow\B,\qquad \chi(q)=0\ \Longleftrightarrow\ q=0,
\]
is a unital semiring homomorphism: a join is zero exactly when all its
summands are zero, and the product assertion is the entirety hypothesis.
For a unital ring $R$, possibly noncommutative or zero, put
\[
 T=Q\times_{\B}R_0.
\]
Thus $T$ has one element over $0\in Q$ and the additive group $R$ over
each nonzero $q$. All transitions between nonzero fibers are identities.
For a left $Q$-semimodule $L$ write
\[
 D=\B\otimes_Q L,\qquad b:L\longrightarrow D,\quad b(l)=1\otimes l.
\]
Here $\B$ is a right $Q$-semimodule through $\chi$. Equivalently, $D$ is
the join-semilattice quotient imposing $ql=l$ for every nonzero $q$.
Every element of $D$ is $b(l)$, and $b(ql)=b(l)$ when $q\ne0$.

\begin{theorem}[Entire support semirings]\label{thm:entire-support-projectives}
The finitely generated projective left $T$-semimodules are classified by
a finitely generated projective left $Q$-semimodule $L$ and a finitely
generated projective left $R$-module $V_p$ for each join-irreducible $p$
of $D=\B\otimes_Q L$. The lattice $D$ is finite and distributive,
although $L$ need not be finite or distributive. Zero labels are allowed.
The corresponding semimodule has support $L$ and groups
\[
 M_l=\bigoplus_{p\le b(l)}V_p,
\]
with the inclusion transitions. A nonzero scalar $(q,r)$ acts by left
multiplication by $r$ into $M_{ql}$, whose displayed summands are the
same as those of $M_l$; the zero scalar acts by the global zero.
Two such data give isomorphic semimodules exactly when a $Q$-isomorphism
of their supports identifies the labels through its induced isomorphism
of Booleanizations.
\end{theorem}
The Booleanization controls the coefficient labels even when the original support is not a distributive lattice. The theorem evaluates this amplitude problem without classifying arbitrary projective $Q$-supports.
